\documentclass[11pt,letterpaper]{article}
\usepackage[margin=1in]{geometry}
\usepackage{amsmath,amssymb,amsthm}
\usepackage{physics}
\usepackage{hyperref}
\usepackage{cleveref}
\usepackage{booktabs}
\usepackage{longtable}
\usepackage{array}

\newtheorem{theorem}{Theorem}[section]
\newtheorem{proposition}[theorem]{Proposition}
\newtheorem{prediction}[theorem]{Prediction}

\numberwithin{equation}{section}
\renewcommand{\theequation}{S12.\arabic{equation}}

\title{Supplement S12: Complete Predictions, Falsification Criteria, and\\
Observational Program \\
\large Supporting Manuscript \S20}
\author{UDT Collaboration}
\date{}

\begin{document}
\maketitle

\begin{abstract}
We compile the complete set of UDT predictions organized by sector:
particle masses, couplings, mixing angles, neutrino properties,
cosmological observables, and laboratory tests. For each prediction we
state the formula, numerical value, current experimental status, and
the falsification criterion. We also outline a program of observational
tests, including searches for undiscovered particles at specific masses,
next-generation neutrino experiments, and cosmological surveys that can
distinguish UDT from $\Lambda$CDM.
\end{abstract}

\tableofcontents

%----------------------------------------------------------------------
\section{Angular sector masses (parameter-free)}
\label{sec:angular_masses}
%----------------------------------------------------------------------

These predictions use only $m_e$ (anchor), $\pi$, and the multiplicities
$(2,3,5,7)$. Zero free parameters.

\begin{longtable}{lccccc}
\toprule
Particle & Formula & Predicted & PDG & Error & Status \\
 & & (MeV) & (MeV) & (\%) & \\
\midrule
\endfirsthead
\toprule
Particle & Formula & Predicted & PDG & Error & Status \\
 & & (MeV) & (MeV) & (\%) & \\
\midrule
\endhead
Pion ($\pi^0$) & $84\pi\,m_e$ & 134.88 & 134.977 & $-0.07$ & Confirmed \\
Muon ($\mu$)   & $\frac{20\pi^3}{3}\,m_e$ & 105.56 & 105.658 & $-0.09$ & Confirmed \\
Proton ($p$)   & $6\pi^5\,m_e$ & 939.07 & 938.272 & $+0.09$ & Confirmed \\
\bottomrule
\end{longtable}

\subsection{Mass ratios (exact)}

\begin{longtable}{lccc}
\toprule
Ratio & Exact form & UDT value & PDG value \\
\midrule
\endfirsthead
$m_p/m_e$ & $6\pi^5$ & 1836.12 & 1836.15 \\
$m_\mu/m_e$ & $20\pi^3/3$ & 206.77 & 206.77 \\
$m_p/m_\pi$ & $\pi^4/14$ & 6.953 & 6.951 \\
$m_\mu/m_\pi$ & $5\pi^2/63$ & 0.783 & 0.783 \\
\bottomrule
\end{longtable}

The tau lepton from Koide $Z_3$ closure (session 2026-03-25):
\begin{equation}
  m_\tau^{\mathrm{Koide}} = 1777.0~\text{MeV}
  \qquad(\text{PDG: }1776.9,\ +0.007\%)\,.
\end{equation}
Brannen parametrization: $A = \sqrt{2j+1} = \sqrt{2}$,
$\theta = (2j+1)/(2\ell+1)^2 = 2/9$.

%----------------------------------------------------------------------
\section{Fine structure constant (three-level chain)}
\label{sec:alpha}
%----------------------------------------------------------------------

\begin{enumerate}
  \item Algebraic bridge: $1/\alpha_0 = 36\pi/I_2 = 137.43$ ($+0.29\%$).
  \item Computed ODE: $\alpha_0 = \mathcal{S}^2 E_1^2 I_2/(2\pi) = 1/137.14$ ($+0.074\%$).
  \item One-loop: $\alpha = \alpha_0(1 + \alpha_0/\pi^2) = 1/137.036$ ($+0.0004\%$).
\end{enumerate}

\noindent\textbf{Critical:} Level~3 uses $\alpha_0$ from Level~2
(computed ODE), not Level~1 (algebraic). See Supplement~S5.

Hydrogen binding: $E_H = \frac{1}{2}\alpha^2 m_e = 13.606$~eV ($-0.0006\%$).

%----------------------------------------------------------------------
\section{Weinberg angle and coupling hierarchy (new)}
\label{sec:weinberg}
%----------------------------------------------------------------------

$\sin^2\theta_W = 3/13 = 0.2308$ (exp: $0.2312$, $-0.19\%$).
$\alpha_s(M_Z) = 0.1178$ (exp: $0.1179$, $-0.07\%$).
See Supplement~S16.

%----------------------------------------------------------------------
\section{CKM matrix (new)}
\label{sec:ckm}
%----------------------------------------------------------------------

$\sin\theta_C = 9/40 = 0.2250$ ($-0.13\%$).
$A = \cos(\pi/5) = 0.8090$ ($-0.24\%$).
$\delta_\mathrm{CKM} = \arctan(9/4) = 66.0^\circ$ ($0.1\sigma$).
$\theta_\mathrm{QCD} = 0$ (exact).
See Supplement~S13.

%----------------------------------------------------------------------
\section{Quark masses and charges (new)}
\label{sec:quarks}
%----------------------------------------------------------------------

$N_c = (2\ell+1) = 3$. $Q_d = -1/3$, $Q_u = +2/3$.
Four inter-generation mass ratios sub-$2\%$:
$m_s/m_d = 2\pi^2$, $m_c/m_u = 6\pi^4$,
$m_b/m_s = 9\pi^2/2$, $m_t/m_c = 14\pi^2$.
See Supplement~S14.

%----------------------------------------------------------------------
\section{Higgs sector (new)}
\label{sec:higgs}
%----------------------------------------------------------------------

$v = 504\pi^6 m_e = 247.6$~GeV ($+0.6\%$).
$\lambda = \pi/24$, $m_H = 126.7$~GeV ($+1.1\%$).
Note: candidate status (not yet derived from metric action).
See Supplement~S15.

%----------------------------------------------------------------------
\section{Nuclear force (new)}
\label{sec:nuclear}
%----------------------------------------------------------------------

$g_A = 4/\pi$ ($-0.23\%$), $f_\pi = 180\,m_e$ ($-0.25\%$),
$g_{\pi NN} = 2\pi^4/15$ ($+0.06\%$),
$g^2/(4\pi) = \pi^7/225$ ($-0.6\%$).
$R_p = \hbar c/(C\varphi_\mathrm{gold}) = 0.865$~fm ($-1.1\%$).
Scalar sector contributes zero ($\alpha_\mathrm{phys} \sim 10^{-40}$).
See Supplement~S16.

%----------------------------------------------------------------------
\section{Proton--neutron mass splitting}
\label{sec:pn}
%----------------------------------------------------------------------

\begin{prediction}
\begin{equation}
  m_n - m_p = \frac{5}{4}\,\alpha_\mathrm{EM}\,C = \frac{5}{4}\,\alpha\,(4\pi^2 m_e r_*).
\end{equation}
\end{prediction}

\noindent\textbf{Predicted:} $\approx 1.28$~MeV.\\
\textbf{PDG:} $1.293$~MeV.\\
\textbf{Error:} $\sim 1\%$.\\
\textbf{Falsification:} Coefficient must be $5/4$ exactly (not $4/3$ or $3/2$).

%----------------------------------------------------------------------
\section{Neutrino sector}
\label{sec:neutrino}
%----------------------------------------------------------------------

\subsection{Neutrino mass}

\begin{prediction}
\begin{equation}
  m_\nu = \frac{\alpha^3 m_e}{(2j+1)^2} = \frac{\alpha^3 m_e}{4}.
\end{equation}
\end{prediction}

\noindent\textbf{Predicted:} $m_\nu \approx 0.050$~eV.\\
\textbf{Current constraint:} $\sqrt{\Delta m^2_{31}} = 0.0501$~eV (NuFIT~5.2, NO).\\
\textbf{Status:} Consistent at $\sim 1\%$.

\subsection{PMNS mixing angles}

\begin{longtable}{lcccc}
\toprule
Angle & UDT & Exact & NuFIT~5.2 & Tension \\
\midrule
\endfirsthead
$\sin^2\theta_{12}$ & 0.30769 & $4/13$ & $0.304 \pm 0.012$ & $0.31\sigma$ \\
$\sin^2\theta_{23}$ & 0.57143 & $4/7$  & $0.573 \pm 0.016$ & $0.10\sigma$ \\
$\sin^2\theta_{13}$ & 0.02222 & $1/45$ & $0.02219 \pm 0.00062$ & $0.05\sigma$ \\
\bottomrule
\end{longtable}

\subsection{CP violation}

\begin{prediction}
UDT predicts $\delta_\mathrm{CP} = 0$ or $\pi$ (no CP violation in the lepton sector).
\end{prediction}

\noindent\textbf{Current status:} NuFIT~5.2 favors $\delta_\mathrm{CP} \approx 197^\circ$
($\sim 1.1\pi$), but the $1\sigma$ range includes $\pi$.\\
\textbf{Falsification:} If $\delta_\mathrm{CP}$ is measured to be outside
$[150^\circ, 210^\circ]$ at $> 5\sigma$, UDT's prediction fails.\\
\textbf{Timeline:} DUNE (2030s), Hyper-Kamiokande (2030s) will measure
$\delta_\mathrm{CP}$ to $\sim 10^\circ$ precision.

\subsection{Mass hierarchy}

\begin{prediction}
UDT predicts normal ordering: $m_1 \approx 0$, $m_2 \approx \sqrt{\Delta m^2_{21}}$,
$m_3 \approx m_\nu = \alpha^3 m_e/4$.
\end{prediction}

\noindent\textbf{Current status:} Normal ordering preferred at $\sim 2.5\sigma$
by NuFIT~5.2.\\
\textbf{Falsification:} Inverted ordering at $> 3\sigma$ would falsify.\\
\textbf{Timeline:} JUNO (2027+), DUNE (2030s).

\subsection{Sum of neutrino masses}

\begin{prediction}
\begin{equation}
  \sum m_\nu \approx m_\nu \approx 0.050~\mathrm{eV}
  \quad (\text{one dominant mass}).
\end{equation}
\end{prediction}

\noindent\textbf{Current constraint:} Cosmological bound $\sum m_\nu < 0.12$~eV
(Planck 2018).\\
\textbf{Falsification:} If $\sum m_\nu > 0.1$~eV is measured, the single
dominant mass prediction fails.\\
\textbf{Timeline:} Euclid, DESI, CMB-S4 (2025--2030).

%----------------------------------------------------------------------
\section{Bridge identities}
\label{sec:bridge}
%----------------------------------------------------------------------

\begin{prediction}
The source--eigenvalue products satisfy:
\begin{align}
  \mathcal{S}^2 E_1^2\big|_{|\kappa|=1} &= 1/18 = 1/(2 \cdot 3^2), \\
  \mathcal{S}^2 E_1^2\big|_{|\kappa|=2} &= \pi^2/36 = \pi^2/(2^2 \cdot 3^2), \\
  \mathcal{S}^2 E_1^2\big|_{|\kappa|=3} &= 1/8 = 1/2^3.
\end{align}
\end{prediction}

\noindent\textbf{Status:} Verified numerically to $< 1\%$ by
\texttt{analysis/04\_bridge\_identities.py}.\\
\textbf{Falsification:} These are exact identities from the metric. A
discrepancy $> 0.1\%$ at machine precision would indicate a computational
error, not a physics failure.

%----------------------------------------------------------------------
\section{Eigenvalue algebraic forms}
\label{sec:eigenvalues}
%----------------------------------------------------------------------

\begin{prediction}
The Dirac ground-state eigenvalues on the vacuum $\phi$ profile take
algebraic values:
\begin{align}
  E_1(\kappa=-1) &= 2\sqrt{2}/3, \\
  E_1(\kappa=-2) &= 16/3, \\
  E_1(\kappa=+2) &= 35/9, \\
  E_1(\kappa=-3) &= 42/5.
\end{align}
\end{prediction}

\noindent\textbf{Status:} Verified to $< 10^{-10}$ by
\texttt{analysis/02\_eigenvalues.py}.\\
\textbf{Falsification:} These are mathematical consequences of the ODE
with the locked parameters. Falsification would require the locked
parameters themselves to be wrong.

%----------------------------------------------------------------------
\section{Cosmological predictions}
\label{sec:cosmology}
%----------------------------------------------------------------------

\subsection{Static universe}

\begin{prediction}
The universe is static. There is no cosmic expansion. The observed redshift
is gravitational: $1 + z = e^{\phi(r)}$.
\end{prediction}

\noindent\textbf{Status:} Currently not falsified (static models can fit
SNe and BAO data; see S10).\\
\textbf{Falsification:} Direct detection of time-varying redshift (Sandage
test) at the predicted $\Lambda$CDM rate $\dot{z} \sim 10^{-10}$~yr$^{-1}$
would falsify static cosmology. UDT predicts $\dot{z} = 0$.\\
\textbf{Timeline:} ELT-ANDES (2030s), 20-year baseline needed.

\subsection{CMB temperature}

\begin{prediction}
\begin{equation}
  T_\mathrm{CMB} = T_\mathrm{starlight}/(1 + z_\mathrm{CMB})
  = 3000/1101 = 2.725~\mathrm{K}.
\end{equation}
\end{prediction}

\noindent\textbf{Status:} COBE FIRAS: $T_\mathrm{CMB} = 2.7255 \pm 0.0006$~K.
Consistent.

\subsection{BAO scale}

\begin{prediction}
The BAO feature is a geometric coherence ruler, not a sound horizon.
The characteristic scale is:
\begin{equation}
  r_d = \frac{1000}{\sqrt{2\pi|k|\,|\beta|\,(1+z_d)}},
\end{equation}
with $k = (3/2)\mu_g$, $\beta = -\cos(\pi/5)\mu_g^2$.
\end{prediction}

\noindent\textbf{Status:} Fits BAO data with $\chi^2/\mathrm{dof} \sim 1.15$.\\
\textbf{Falsification:} If the BAO feature is shown to depend on baryon
density $\Omega_b$ (not just geometry), the geometric interpretation fails.

\subsection{Self-consistency}

\begin{prediction}
\begin{equation}
  c^2 = \frac{2GM}{r_*}
\end{equation}
(the universe is at the Misner--Sharp horizon).
\end{prediction}

\noindent\textbf{Status:} Verified to $0.0000\%$ with BAO-fitted parameters.

%----------------------------------------------------------------------
\section{Laboratory tests}
\label{sec:laboratory}
%----------------------------------------------------------------------

\subsection{Precision mass ratios}

\begin{prediction}
$m_p/m_e = 6\pi^5 = 1836.118\ldots$
\end{prediction}

\noindent\textbf{PDG:} $m_p/m_e = 1836.15267\ldots$\\
\textbf{Discrepancy:} $0.019\%$ (34~ppm).\\
\textbf{Test:} Improved measurement precision (currently $\sim 0.1$~ppb)
tests whether the deviation is real or whether UDT needs a small
higher-order correction.

\subsection{Muon anomalous magnetic moment}

UDT predicts corrections to $g-2$ from the vacuum $\phi$ profile. The
anomalous magnetic moment should receive a geometric correction of order:
\begin{equation}
  \delta(g-2)_\mu \sim \frac{\alpha}{\pi}\left(1 + \mathcal{O}(e^{2\phi_0})\right).
\end{equation}

The leading term matches the Schwinger result; subleading terms may differ
from QED at the $\sim 10^{-10}$ level.

%----------------------------------------------------------------------
\section{Undiscovered particles}
\label{sec:undiscovered}
%----------------------------------------------------------------------

\subsection{Predictions from the radial eigenvalue spectrum}

Higher Dirac eigenvalues $E_n(\kappa)$ predict specific meson and baryon
masses. States not yet observed (or poorly measured) provide discovery
targets:

\begin{longtable}{lcccc}
\toprule
Channel & $n$ & $E_n$ (predicted) & Mass (MeV) & Notes \\
\midrule
\endfirsthead
$\kappa=+1$ & 3 & (computed) & $\sim 1400$ & Unassigned meson \\
$\kappa=-2$ & 2 & (computed) & $\sim 1800$ & f2-family excitation \\
$\kappa=+3$ & 1 & (computed) & $\sim 1100$ & Exotic quantum numbers \\
$\kappa=-3$ & 2 & (computed) & $\sim 2400$ & High-mass meson \\
\midrule
\multicolumn{5}{l}{\textit{Baryon sub-cavity ($r_b = r_*/\varphi$):}} \\
$\kappa=-2$ & 2 & (computed) & $\sim 1800$ & Baryon excitation \\
$\kappa=+2$ & 2 & (computed) & $\sim 2200$ & $\Omega$-family member \\
\bottomrule
\end{longtable}

Exact values require running \texttt{analysis/02\_eigenvalues.py} with
extended mode counts.

\subsection{The $\kappa = +3$ sector}

The $\kappa = +3$ channel has no known algebraic eigenvalue form. Its
ground-state eigenvalue predicts a particle mass in the $\sim 1$--$1.2$~GeV
range. If this state exists and has the predicted mass, it would be strong
evidence for the $|\kappa_\mathrm{max}| = 3$ selection.

%----------------------------------------------------------------------
\section{Falsification criteria: summary table}
\label{sec:falsification}
%----------------------------------------------------------------------

\begin{longtable}{p{4cm}p{4cm}p{4cm}}
\toprule
Prediction & Falsification criterion & Timeline \\
\midrule
\endfirsthead
\toprule
Prediction & Falsification criterion & Timeline \\
\midrule
\endhead
$m_\pi = 84\pi m_e$ & $> 0.5\%$ deviation from exact formula & Already testable \\
$m_p = 6\pi^5 m_e$ & $> 0.5\%$ deviation & Already testable \\
$\sin^2\theta_{12} = 4/13$ & Measured outside $[0.28, 0.33]$ at $5\sigma$ & JUNO (2027+) \\
$\sin^2\theta_{23} = 4/7$ & Measured outside $[0.53, 0.61]$ at $5\sigma$ & DUNE (2030+) \\
$\sin^2\theta_{13} = 1/45$ & Measured outside $[0.018, 0.026]$ at $5\sigma$ & Current data \\
$\delta_\mathrm{CP} = 0$ or $\pi$ & $|\delta_\mathrm{CP} - \pi| > 30^\circ$ at $5\sigma$ & DUNE/HK (2030+) \\
$m_\nu = \alpha^3 m_e/4$ & $m_\nu > 0.1$~eV or inverted ordering at $5\sigma$ & JUNO/DUNE \\
Normal mass ordering & Inverted ordering at $> 3\sigma$ & JUNO (2027+) \\
$\sum m_\nu \approx 0.05$~eV & $\sum m_\nu > 0.12$~eV at $3\sigma$ & Euclid/DESI \\
Static universe ($\dot{z}=0$) & $\dot{z} \neq 0$ at $3\sigma$ & ELT (2040+) \\
$c^2 = 2GM/r_*$ & $> 1\%$ violation & (Theoretical) \\
Bridge identities & $> 0.1\%$ violation at machine precision & (Computational) \\
Three generations only & Fourth neutrino detected & Short-baseline \\
\bottomrule
\end{longtable}

%----------------------------------------------------------------------
\section{Ancient remnants survey program}
\label{sec:remnants}
%----------------------------------------------------------------------

\subsection{Motivation}

If UDT is correct, the universe is static and has no finite age. This
implies that ``ancient remnants'' -- structures formed at arbitrarily early
times -- may exist. The detection or non-detection of such remnants provides
a test of the static cosmology.

\subsection{Observable signatures}

\begin{enumerate}
  \item \textbf{Ultra-high-redshift galaxies:} In a static universe,
    galaxies at $z > 15$ should exist with fully formed stellar populations.
    JWST has already detected surprisingly mature galaxies at $z \sim 13$.
    UDT predicts these as ordinary galaxies at large $r$ (high gravitational
    redshift), not ``impossibly early'' objects.

  \item \textbf{Metal-enriched intergalactic medium at all $z$:} In $\Lambda$CDM,
    the IGM should be metal-poor at $z > 6$ (before significant star formation).
    In UDT, metals have had unlimited time to disperse.

  \item \textbf{Gravitational wave background:} A static universe accumulates
    gravitational waves from all stellar events at all times, producing a
    stochastic background with specific spectral characteristics.

  \item \textbf{CMB spectral distortions:} The static-universe CMB is a
    thermal equilibrium state, not a relic of a hot phase. Any spectral
    distortion ($\mu$, $y$) must arise from foreground effects, not primordial
    processes. PIXIE-class experiments can test this.

  \item \textbf{Nucleosynthesis ratios:} In a static universe, light element
    abundances (H/He/Li) must be produced by stellar processes or maintained
    by an equilibrium mechanism, not by Big Bang nucleosynthesis. The
    observed $^4$He mass fraction $Y_p \approx 0.245$ must be explained
    without invoking a hot early universe.
\end{enumerate}

\subsection{Proposed observational program}

\begin{longtable}{p{4cm}p{4cm}p{3.5cm}}
\toprule
Test & Facility & Timeline \\
\midrule
\endfirsthead
$z > 15$ galaxies & JWST, Roman Space Telescope & 2025--2030 \\
IGM metallicity at $z > 6$ & ELT, TMT & 2030--2035 \\
Sandage test ($\dot{z}$) & ELT-ANDES & 2035--2055 \\
CMB distortions & PIXIE, Voyage 2050 & 2030+ \\
Stochastic GW background & LISA, Einstein Telescope & 2035+ \\
Primordial He abundance & Ground-based surveys & Ongoing \\
\bottomrule
\end{longtable}

%----------------------------------------------------------------------
\section{Hierarchy of predictions by confidence}
\label{sec:hierarchy}
%----------------------------------------------------------------------

We classify predictions by their derivation status:

\subsection{Tier 1: Exact algebraic (highest confidence)}

These follow directly from the metric and Diophantine identity with no
numerical computation:
\begin{itemize}
  \item $m_\pi/m_e = 84\pi$
  \item $m_\mu/m_e = 20\pi^3/3$
  \item $m_p/m_e = 6\pi^5$
  \item $\sin^2\theta_{12} = 4/13$, $\sin^2\theta_{23} = 4/7$,
    $\sin^2\theta_{13} = 1/45$
  \item Bridge identities: $1/18$, $\pi^2/36$, $1/8$
  \item Eigenvalue algebraic forms
\end{itemize}

\subsection{Tier 2: Numerical from the metric (high confidence)}

These require solving the ODE but have no free parameters:
\begin{itemize}
  \item $1/\alpha = 36\pi/I_2$
  \item Individual eigenvalues $E_n(\kappa)$ and mass predictions
  \item $m_n - m_p = (5/4)\alpha C$
  \item $m_\nu = \alpha^3 m_e/4$
\end{itemize}

\subsection{Tier 3: Structural (moderate confidence)}

These follow from the framework but involve interpretive steps:
\begin{itemize}
  \item $\delta_\mathrm{CP} = 0$ or $\pi$
  \item Three generations (no fourth)
  \item Normal mass ordering
  \item Static cosmology
\end{itemize}

\subsection{Tier 4: Extended (lower confidence)}

These are consequences of the static cosmology that depend on the
cosmological profile:
\begin{itemize}
  \item BAO geometric ruler
  \item $\dot{z} = 0$
  \item Nucleosynthesis mechanism
  \item Ancient remnants
\end{itemize}

%----------------------------------------------------------------------
\section{Verification}
\label{sec:verification}
%----------------------------------------------------------------------

The full prediction table draws on all analysis scripts in the repository
(01--10) and all supplements (S1--S11). Each prediction is tagged with its
verification script and output file. The complete pipeline can be re-run
with:
\begin{verbatim}
  cd /home/udt-admin/udt_repo
  for f in analysis/0*.py; do python3 $f; done
\end{verbatim}

All predictions should remain within their stated error bounds. Any
deviation indicates either a computational error (to be debugged) or a
genuine falsification (to be reported).

\end{document}
